\documentclass[12pt,a4paper]{article}
\usepackage[margin=25mm, headheight=15pt, headsep=10pt]{geometry}
\usepackage{fontspec}
\usepackage[table]{xcolor} % Note: table option is required for rowcolors
\usepackage{titlesec}
\usepackage{tocloft}
\usepackage{fancyhdr}
\usepackage{booktabs}
\usepackage{tcolorbox}
\tcbuselibrary{skins, breakable}
\usepackage[colorlinks=true, linkcolor=emerald, urlcolor=emerald, bookmarks=true]{hyperref}

\definecolor{emerald}{RGB}{0,102,51}
\definecolor{darkred}{RGB}{139,0,0}
\definecolor{lightgray}{RGB}{245,245,245}

\usepackage[nil]{babel}
\babelprovide[import, main]{bengali}
\babelprovide[import]{arabic}
\babelfont{rm}[Renderer=HarfBuzz,Scale=1.0]{Noto Serif Bengali}
\babelfont{sf}[Renderer=HarfBuzz]{Noto Sans Bengali}
\babelfont{tt}[Renderer=HarfBuzz]{Noto Sans Bengali}
\babelfont[arabic]{rm}[Renderer=HarfBuzz,Scale=1.1]{Noto Naskh Arabic}

% Section styling
\titleformat{\section}{\Large\bfseries\color{emerald}}{\thesection.}{0.5em}{}
\titleformat{\subsection}{\large\bfseries\color{emerald!85!black}}{\thesubsection.}{0.5em}{}
\titleformat{\subsubsection}{\normalsize\bfseries\color{emerald!70!black}}{\thesubsubsection.}{0.5em}{}
\titlespacing*{\section}{0pt}{18pt}{8pt}

% Fancy Header/Footer setup
\pagestyle{fancy}
\fancyhf{} % clear all header and footer fields
\fancyhead[L]{\color{emerald}\small\textbf{\nouppercase{\leftmark}}}
\fancyfoot[L]{\scriptsize\color{black!50}{\lat{AI}}-সংকলিত, আলেম-যাচাইকৃত নয়}
\fancyfoot[C]{\color{emerald!70!black}\thepage}
\renewcommand{\headrulewidth}{0.4pt}
\renewcommand{\headrule}{\hbox to\headwidth{\color{emerald}\leaders\hrule height \headrulewidth\hfill}}
\renewcommand{\footrulewidth}{0pt}

% Custom boxes
% 1. Ayat/Hadith Box
\newtcolorbox{ayatbox}[1][]{
    enhanced, breakable, 
    colback=emerald!5, 
    colframe=emerald, 
    boxrule=0.5pt, 
    arc=3pt, 
    left=10pt, right=10pt, top=10pt, bottom=10pt,
    #1
}

% 2. Quote/Translation Box
\newtcolorbox{quotebox}[1][]{
    enhanced, breakable,
    frame hidden,
    colback=lightgray,
    borderline west={3pt}{0pt}{emerald},
    left=15pt, right=10pt, top=10pt, bottom=10pt,
    sharp corners,
    #1
}

% 3. AI-generated notice box
\newtcolorbox{warnbox}[1][]{
    enhanced, breakable,
    colback=red!4,
    colframe=darkred,
    boxrule=0.8pt,
    arc=3pt,
    left=10pt, right=10pt, top=10pt, bottom=10pt,
    #1
}

% Commands
\newcommand{\arab}[1]{\foreignlanguage{arabic}{#1}}
\newcommand{\kib}[1]{\textcolor{emerald}{\textbf{#1}}}

% Latin (English) fallback: Noto Serif Bengali has NO Latin glyphs
\newfontfamily\latfont[Renderer=HarfBuzz]{Noto Serif}
\newcommand{\lat}[1]{{\latfont #1}}

\setlength{\parskip}{5pt}
\setlength{\parindent}{0pt}

% Fix for missing bullet in Noto Serif Bengali
\renewcommand{\labelitemi}{$\bullet$}



\begin{document}
\begin{center}{\Large\bfseries\color{emerald} \lat{01}-তাকবীরে-তাহরিমা}\end{center}
\vspace{1em}
\section*{০১ — তাকবীরে তাহরিমা: \textit{\lat{Allahu} \lat{Akbar}} — নামাজের প্রথম কথা}

\begin{warnbox}
 \textbf{গুরুত্বপূর্ণ নোটিশ (\lat{Important} \lat{Notice}):} এই কন্টেন্টটি \lat{AI} (কৃত্রিম বুদ্ধিমত্তা) দিয়ে প্রস্তুত ও সংকলিত। \lat{AI} কঠোর সূত্র-মান নীতি মেনে শুধু নির্ভরযোগ্য উৎস থেকে তথ্য সংগ্রহ করেছে — কেবল সহীহ/হাসান হাদিস ও সহীহ সনদের আসার রাখা হয়েছে, দুর্বল সনদ বাদ দেওয়া হয়েছে। তবে এটি কোনো আলেম বা মুহাদ্দিস কর্তৃক যাচাইকৃত নয়।
\end{warnbox}

\medskip\hrule\medskip

\subsection*{১. সূত্র কার্ড}

\begin{table}[h]\centering\small
\begin{tabular}{ll}
বিষয় & বিবরণ \\
\hline
\textbf{তাকবীর} & \textbf{\arab{اللَّهُ} \arab{أَكْبَرُ}} (\textit{\lat{Allahu} \lat{Akbar}}) — আল্লাহ সবচেয়ে বড় \\
\textbf{বিষয়} & নামাজের সূচনা; রুকন; সকল তাকবীরের অর্থ \\
\textbf{সূত্র} & সহীহ বুখারি ৭৩৫, সহীহ মুসলিম ৩৯০ \\
\textbf{প্রেক্ষাপট} & কিবলামুখী হয়ে দাঁড়িয়ে, হাত উঁচু করে, "আল্লাহু আকবার" বলে নামাজ শুরু করা \\
\hline
\end{tabular}
\end{table}

\medskip\hrule\medskip

\subsection*{২. মূল আরবি}

\begin{center}\arab{اللَّهُ أَكْبَرُ}\end{center}

\medskip\hrule\medskip

\subsection*{৩. বাংলা উচ্চারণ}

\textbf{আল্লাহু আকবার}

\medskip\hrule\medskip

\subsection*{৪. শব্দ-শব্দ অর্থ}

\begin{table}[h]\centering\small
\begin{tabular}{lll}
আরবি & বাংলা & \lat{English} \\
\hline
\textbf{\arab{اللَّهُ}} & আল্লাহ & \textit{\lat{Allah}} \\
\textbf{\arab{أَكْبَرُ}} & সবচেয়ে বড় / সর্বশ্রেষ্ঠ & \textit{\lat{the} \lat{Greatest}} \\
\hline
\end{tabular}
\end{table}

\medskip\hrule\medskip

\subsection*{৫. পূর্ণ বাংলা অনুবাদ}

\textbf{\lat{“}আল্লাহ সবচেয়ে বড়।\lat{”}}

\medskip\hrule\medskip

\subsection*{৬. সংক্ষিপ্ত ব্যাখ্যা (\lat{Khushu-focused})}

\subsubsection*{\textbf{তাকবীরে তাহরিমা কী?}}
\textbf{\lat{Takbirat} \lat{al-Ihram} (তাকবীরে তাহরিমা)} হলো নামাজের প্রথম তাকবীর — এটি দিয়ে নামাজ \textbf{\lat{haram} (নিষিদ্ধ)} অবস্থায় প্রবেশ করা হয় (যার অর্থ: এরপর নামাজের বাইরে কিছু করা যায় না — হাত নড়ানো, কথা বলা, খাওয়া-দাওয়া)। এটি নামাজের \textbf{\lat{pillar} (রুকন)} — ছাড়লে নামাজ হয় না।

\subsubsection*{\textbf{"আকবার" — কী বড়?}}
\textbf{"\lat{Akbar} (আকবার)"} মানে "সবচেয়ে বড়" — সব কিছু থেকে বড়। এই শব্দে আছে \textbf{\lat{tawhid} (তাওহীদ)}-এর সার: আমি যা কিছু ভয় পাই, ভালোবাসি, আশা করি, ভরসা করি — সব আল্লাহর চেয়ে \textbf{\lat{lesser} (কম)}। তাই নামাজে প্রবেশ করেই প্রথমেই এই \textbf{\lat{acknowledgment} (স্বীকৃতি)} দিচ্ছি।

\subsubsection*{\textbf{কীভাবে বলবেন?}}
\begin{itemize}
\item \textbf{ধীরে ধীরে}, \textbf{স্পষ্ট করে}, \textbf{শ্রদ্ধায়}।
\item \textbf{উচ্চ স্বরে} (নিজে শোনার মতো, অন্যকে নয়)।
\item \textbf{হাত উঁচু করুন} — কানের সমান বা কান্ধের সমান; \textbf{ধীরে} উঁচু করুন, \textbf{ধীরে} নামান।
\item \textbf{কানে স্পর্শ করুন না} — কানের সমানে আনুন মাত্র।
\end{itemize}
\subsubsection*{\textbf{কখন বলবেন?}}
\begin{itemize}
\item কিবলামুখী হয়ে দাঁড়িয়ে, \textbf{নিয়ত} করে (মনে), হাত উঁচু করে।
\item এরপর \textbf{বিসমিল্লাহ} বলে সানা শুরু করুন।
\end{itemize}
\medskip\hrule\medskip

\subsection*{৭. নামাজের সকল তাকবীরের তালিকা}

\begin{table}[h]\centering\small
\begin{tabular}{lll}
\# & তাকবীর & কখন বলা হয় \\
\hline
১ & \textbf{\arab{تَكْبِيرَةُ} \arab{التَّحْرِيمِ}} — "আল্লাহু আকবার" (প্রথম) & নামাজ শুরু করার সময় \\
২ & রুকুতে যাওয়ার সময় & "আল্লাহু আকবার" \\
৩ & রুকু থেকে উঠে দাঁড়ানোর সময় & "সামিআল্লাহু লিমান হামিদাহ" \\
৪ & সিজদায় যাওয়ার সময় & "আল্লাহু আকবার" \\
৫ & সিজদা থেকে উঠে বসার সময় & "আল্লাহু আকবার" \\
৬ & দ্বিতীয় সিজদা থেকে উঠে দাঁড়ানোর সময় & "আল্লাহু আকবার" \\
৭ & তৃতীয়/চতুর্থ রাকাতে কিয়ামে উঠার সময় & "আল্লাহু আকবার" \\
৮ & সালামের আগে (ইমাম) & "আল্লাহু আকবার" \\
\hline
\end{tabular}
\end{table}

\textbf{মোট:} প্রতিটি রাকাতে প্রায় \textbf{৭-৮ বার} "আল্লাহু আকবার" বলা হয়।

\medskip\hrule\medskip

\subsection*{৮. সংশ্লিষ্ট সহীহ হাদিস}

\begin{table}[h]\centering\small
\begin{tabular}{lll}
বিষয় & হাদিস & গ্রন্থ ও নম্বর \\
\hline
নামাজ কীভাবে পড়বেন & \textbf{\lat{“}নামাজ পড়ো, যেমন আমাকে নামাজ পড়তে দেখেছ।\lat{”}} & সহীহ বুখারি ৬৩১ \\
তাকবীরে তাহরিমা রুকন & \textbf{\lat{“}তোমাদের প্রত্যেকে নিজের ইমামের কথা মেনে চলুক।\lat{”}} & সহীহ বুখারি ৭৩৫ \\
তাকবীরে তাহরিমার বৈশিষ্ট্য & হাত উঁচু করা সুন্নাত & সুনানে আবু দাউদ ২৫৯০ (সহীহ) \\
\hline
\end{tabular}
\end{table}

\medskip\hrule\medskip

\subsection*{৯. খুশুর জন্য মূল পয়েন্টস}

\begin{table}[h]\centering\small
\begin{tabular}{ll}
পয়েন্ট & কীভাবে প্রয়োগ করবেন \\
\hline
\textbf{"আকবার" শব্দের গভীর অর্থ} & হাত উঁচু করার সময় মনে করুন: \textbf{"আমি যা কিছু ভয় পাই, সব আল্লাহর চেয়ে ছোট।"} \\
\textbf{দুনিয়া ছাড়া} & তাকবীর বলার পর আপনি \textbf{\lat{dunya} (দুনিয়া)} থেকে \textbf{\lat{bari} (বিরত)} হলেন — এটি \textbf{\lat{symbol} (প্রতীক)}। \\
\textbf{স্পষ্ট ও ধীর} & ধীরে বলুন — \textbf{\lat{moment} (মুহূর্ত)} টেনে নিন। \\
\textbf{হাতের গতি} & হাত \textbf{ধীরে} উঁচু করুন, \textbf{ধীরে} নামান। \textbf{\lat{Hurried} (তাড়াহুড়ো)} করবেন না। \\
\textbf{কিবলার দিকে তাকিয়ে} & চোখ কিবলার দিকে স্থির — \textbf{\lat{focus} (মনোযোগ)} ভাঙবেন না। \\
\hline
\end{tabular}
\end{table}

\medskip\hrule\medskip

\textit{শেষ আপডেট: আগস্ট ২০২৬}

\end{document}
